\section{Обсуждение и ограничения}
Полное время включает постоянные и переменные составляющие. Условная модель
\[
 T_f(N,q)=A_f+B_f(N,q)+\varepsilon
\]
помогает интерпретировать результаты: $A_f$ включает создание процесса и инфраструктуру, $B_f$ -- работу с входом, а $\varepsilon$ -- изменчивость среды. По трём точкам нельзя надёжно оценить универсальную функцию затрат. Если время почти не изменилось при росте небольшого входа, это согласуется с существенной постоянной составляющей, но само по себе не определяет её точный размер.

Низкая средняя CPU-загрузка не обязательно означает меньшее процессорное время. Процесс может долго ожидать файловую систему и иметь малую долю CPU при большом времени ответа. Поэтому $\widehat C$ и $\widehat U$ показаны раздельно. При сравнении ресурсов приоритет имеет накопленное время, а процент помогает объяснить соотношение работы и ожидания.

Один worker Vitest не означает один поток операционной системы. Node.js и вспомогательные компоненты могут использовать собственные потоки; часть служб может существовать как дочерний процесс. Ограничивается параллелизм тестовых задач, а не всё внутреннее устройство среды. Сравнение не пытается сделать реализации раннеров идентичными: различия их инфраструктуры входят в предмет исследования.

Наблюдатель с дискретным опросом не гарантирует фиксацию короткого пика памяти. Отсчёты RSS разных процессов также не являются строго одновременными. Записанный максимум характеризует используемый метод измерения. Для точного учёта завершённых процессов потребовались бы механизмы ОС, например специализированная трассировка и учёт дерева задач, которые выходят за рамки компактного кроссплатформенного стенда.

Наконец, установка и прогрев могут менять состояние кэшей. Измерения выполняются на рабочей Windows-машине, а не на выделенном лабораторном сервере. В отчёте не заявлено, что все фоновые процессы были отключены. Повторный запуск способен дать другие абсолютные величины. Сохранённые версии и входы обеспечивают повторяемую процедуру, но не обещают побитовое совпадение временных показателей.

\section{Инструкция по установке и запуску}
Для программы требуются Node.js версии не ниже 22 с npm и Python версии не ниже 3.10. В данной апробации использованы конкретные версии, указанные в разделе результатов. Команды выполняются в каталоге \code{benchmark_app}. Интернет нужен при первой установке пакетов; измеряемые тесты после установки сетевых обращений не требуют.

\begin{lstlisting}
cd benchmark_app
npm ci
python -m venv .venv
.venv\Scripts\python.exe -m pip install -r requirements.txt
npm run verify
npm test
npm run demo
npm start -- --text "abba"
npm start -- --file data/input-example.json
npm run test:jest
npm run test:mocha
npm run test:vitest
\end{lstlisting}

Команды выше рассчитаны на Windows PowerShell. Активация окружения не обязательна: можно непосредственно вызывать его Python. В Linux и macOS вместо \code{.venv\Scripts\python.exe} используют \code{.venv/bin/python}. Программа JavaScript не требует отдельной компиляции. Команда \code{npm ci} устанавливает версии из lock-файла и пригодна для повторяемого code review.

Заготовленные входы включены в проект. Для пересоздания или изменения размера:
\begin{lstlisting}
.venv\Scripts\python.exe generate.py
.venv\Scripts\python.exe generate.py --sizes 200 2000 --cases 12
node src/demo.js data/author-palindromes.csv
\end{lstlisting}
Увеличение \code{--cases} повышает число записей, а \code{--sizes} -- длину каждой строки. После изменения нужно заново измерить производительность; старые результаты не описывают новые файлы.

Короткая проверка измерителя:
\begin{lstlisting}
.venv\Scripts\python.exe benchmark.py --input data/size-100.json --runs 2 --output reports/smoke
\end{lstlisting}
Основная серия:
\begin{lstlisting}
.venv\Scripts\python.exe benchmark.py --runs 7 --repeats 10 --output reports/new-json
\end{lstlisting}
Эквивалентная серия CSV:
\begin{lstlisting}
.venv\Scripts\python.exe benchmark.py --input data/size-1000.csv --runs 7 --output reports/new-csv
\end{lstlisting}

Если Node.js не находится в PATH, путь к исполняемому файлу передаётся через \code{--node}. Параметр \code{--timeout} задаёт предел в секундах. \code{--interval-ms} задаёт паузу опроса; изменение этого параметра требует новой серии и должно отражаться в интерпретации результатов. Запись производится в каталог \code{--output}, поэтому для сохранения прошлой серии выбирают новое имя.

Полная автоматическая приёмка запускается командой \code{.venv\Scripts\python.exe validate.py}. Она проверяет все три фреймворка, включая намеренно неверное ожидание. Сохранённые результаты новой апробации находятся в \code{reports/2026-09-16-json} и \code{reports/2026-09-16-csv}. Существующая серия не перезаписывается: при повторе следует выбрать новое имя каталога.

Для сборки документов нужны XeLaTeX и Biber из TeX Live или MiKTeX, пакеты Beamer, biblatex, algorithm и algpseudocode, а также шрифты Times New Roman, Arial и Consolas. Из корня проекта:
\begin{lstlisting}
xelatex -interaction=nonstopmode -halt-on-error My_thesis.tex
biber My_thesis
xelatex -interaction=nonstopmode -halt-on-error My_thesis.tex
xelatex -interaction=nonstopmode -halt-on-error My_thesis.tex
xelatex -interaction=nonstopmode -halt-on-error Presentation.tex
biber Presentation
xelatex -interaction=nonstopmode -halt-on-error Presentation.tex
xelatex -interaction=nonstopmode -halt-on-error Presentation.tex
\end{lstlisting}
Готовые PDF приложены, поэтому установка LaTeX не нужна для чтения и демонстрации. Для смены персональных данных редактируются титульные блоки двух исходников. При пересчёте экспериментальных результатов сначала обновляются автоматически генерируемые таблицы и диаграммы согласно README.

\section{Заключение}
По принятому правилу для исследуемого приложения выбран Mocha. Он получил 9 первых мест из девяти сравнений. При длине нагрузочной строки 10000 его медианы составили: 601.0 мс полного времени, 484.4 мс CPU-времени и 101.1 MiB наблюдаемого пика RSS. Выбор относится к полному набору из 99 тестов на данной машине; он не устанавливает универсального победителя для всех проектов.


Разработан компактный стенд сравнения Jest, Mocha и Vitest на общей алгоритмической нагрузке. Он импортирует файлы переменного размера в JSON и CSV, проверяет результат по эталону, выполняет одинаковые утверждения в трёх раннерах и сохраняет данные для воспроизведения эксперимента.

Содержательная проверка предшествует измерению скорости. Независимый перебор коротких строк, специальные случаи, метаморфное свойство и негативные входы проверяют разные классы ошибок. Пошаговый пример показывает связь между радиусами Манакера, найденной подстрокой и количеством палиндромных вхождений.

Сравнение учитывает полное время запуска, накопленное CPU-время и наблюдаемую память дерева процессов. Определения метрик позволяют не смешивать CPU-проценты с процессорным временем, а RSS -- с размером JavaScript-кучи. Ограничения дискретного опроса явно включены в методику и трактовку результатов.

Практический результат состоит в проверяемой процедуре выбора для заданной конфигурации. Полученные числа применимы к одному пакетному тесту чистой функции и указанной машине. Следующее развитие исследования -- несколько тестовых файлов, режим наблюдения, покрытие и прикладной набор проекта. Эти сценарии требуют отдельных серий, а не распространения текущего ранжирования без измерений.
\label{body-end}
